
% =============================================================================
% Table 02 — Defining Invariants
% Book I — Ontology of Viability
% Chapter 02 — Invariance
%
% Purpose:
%   To identify the analytical conditions and evaluative criteria
%   required for establishing an invariant under transformation.
%
% Conceptual Policy:
%   1. A formal invariant is not automatically an ontological invariant.
%   2. Preservation must be specified relative to a transformation.
%   3. The domain and respect of invariance must be explicit.
%   4. Invariance is not a criterion or definition of identity.
%   5. The listed conditions are analytical requirements, not
%      independent ontological entities.
% =============================================================================

\begin{table}[htbp]
    \centering
    \small
    \renewcommand{\arraystretch}{1.4}
    \setlength{\tabcolsep}{5pt}

    \caption{Analytical Conditions for Identifying Invariants}
    \label{tab:defining-invariants}

    \begin{tabular}{
        p{2.7cm}
        p{3.4cm}
        p{4.0cm}
        p{3.1cm}
    }
        \hline
        \textbf{Condition}
        &
        \textbf{Central Question}
        &
        \textbf{Analytical Requirement}
        &
        \textbf{Risk if Unspecified}
        \\
        \hline

        Subject of Invariance
        &
        What is claimed to remain invariant?
        &
        Identify the existent, determination, relation, magnitude,
        or configuration under consideration.
        &
        Attribution of invariance to an unspecified object or feature.
        \\
        \hline

        Transformation
        &
        Under what transformation is preservation asserted?
        &
        Specify the relevant change, operation, or transformation
        against which preservation is assessed.
        &
        Treating invariance as an unrestricted or absolute property.
        \\
        \hline

        Domain of Transformation
        &
        Within which range of conditions does preservation obtain?
        &
        Delimit the states, configurations, conditions, or circumstances
        over which the transformation is considered.
        &
        Illegitimate generalization beyond the conditions examined.
        \\
        \hline

        Respect of Invariance
        &
        In what particular respect does preservation occur?
        &
        Specify the determination, relation, magnitude, or organizational
        feature that remains preserved.
        &
        Confusing partial preservation with total preservation.
        \\
        \hline

        Criterion of Preservation
        &
        What establishes that the relevant feature remains preserved?
        &
        Identify the appropriate criterion of sameness or preservation
        for the specified determination or relation.
        &
        Confusing formal resemblance, approximation, or descriptive
        similarity with actual preservation.
        \\
        \hline

        Transformation Admissibility
        &
        Which transformations are relevant to the claim?
        &
        Distinguish transformations included within the relevant
        domain from those excluded by the claim's conditions.
        &
        Arbitrary selection of transformations to produce an apparent
        invariant.
        \\
        \hline

        Ontological Grounding
        &
        What is the ontological status of the preserved feature?
        &
        Establish whether the preserved feature corresponds to an
        actual determination or relation of the existent, rather
        than merely a formal feature of its representation.
        &
        Equating formal invariance with ontological preservation.
        \\
        \hline

        Scope of Validity
        &
        How far does the invariance claim extend?
        &
        State whether preservation is conditional, restricted,
        or established throughout the specified domain.
        &
        Extending a limited invariance claim to unrestricted
        transformations or circumstances.
        \\
        \hline

    \end{tabular}

    \vspace{0.6em}

    \begin{minipage}{0.95\textwidth}
        \footnotesize
        \textit{Note:}
        The conditions listed above provide an analytical framework
        for evaluating claims of invariance. They do not constitute
        a universally sufficient test for the ontological existence
        of an invariant. Their relevance and adequacy may vary
        according to the nature of the determination, relation,
        and transformation under consideration.

        The identification of a formally preserved quantity or
        relation does not, by itself, establish the numerical identity,
        persistence, or continuity of the existent as a whole.
    \end{minipage}

\end{table}
